\section{A uniform bound for the complete collar flux}
\label{sec:upper-bound}

The spatial estimate in this section holds for arbitrary real frequencies.
The passage to a heat-flow estimate uses the additional assumption that
the frequencies lie in one dyadic band.  The proof of the spatial estimate
follows the half-plane and Laguerre-polynomial approach of
Zhang~\cite{Zhang2026}.  The details, including the finite-range step, are
given below.  The constants are chosen for a short proof.

\subsection{A finite-dimensional half-plane basis}

For distinct complex numbers $\zeta_1,\ldots,\zeta_N$ with
$b_j=\operatorname{Re}\zeta_j>0$, consider the space spanned by
$e^{-\zeta_jt}$ in $L^2(0,\infty)$.  The Laplace transform, with boundary
norm $(2\pi)^{-1}\int_{\mathbb R}|\widehat h(i\xi)|^2\,d\xi$, is an
isometry on this space by Fourier Plancherel.  A useful orthonormal basis
has transforms
\begin{equation}\label{eq:ub-basis}
 \widehat\varphi_j(s)=\frac{\sqrt{2b_j}}{s+\zeta_j}
       \prod_{r<j}\frac{s-\overline{\zeta_r}}{s+\zeta_r},
 \qquad 1\le j\le N.
\end{equation}
Here and below inner products are linear in the first argument.
To verify this assertion, note that each factor in the product has
modulus one on the imaginary axis.  The kernel
$k_w(s)=(s+\overline w)^{-1}$ satisfies
\[
 \langle H,k_w\rangle=H(w),\qquad
 \|k_w\|^2=(2\operatorname{Re}w)^{-1}.
\]
These identities follow by applying Plancherel to
$e^{-\overline w t}$.  The norm of every vector in
\eqref{eq:ub-basis} is therefore one.  In the inner product of the $j$th
and $\ell$th vectors, $\ell<j$, cancel their common initial factors.
The remaining $j$th factor vanishes at $\overline{\zeta_\ell}$, so the
reproducing identity gives orthogonality.  Partial fractions show that
the first $j$ vectors lie in the span of the first $j$ kernels.  The
coefficient at $(s+\zeta_j)^{-1}$ is nonzero, since
$\zeta_r-\zeta_j\ne0$ and $\operatorname{Re}(\zeta_j+
\overline{\zeta_r})>0$.  This proves spanning as well.

Let
\[
 (V_\zeta h)(t)=\int_0^t e^{-\zeta(t-u)}h(u)\,du.
\]
Inverting the factors in \eqref{eq:ub-basis} gives
\begin{equation}\label{eq:ub-volterra}
 \varphi_j=
 \prod_{r<j}(I-2b_rV_{\zeta_r})
       \bigl(\sqrt{2b_j}e^{-\zeta_jt}\bigr).
\end{equation}
These are identities between entire exponential sums, so the same formula
holds at negative real times, with each integral oriented from $0$ to $t$.

Write
\[
 L_m(-y)=\sum_{\ell=0}^m\binom m\ell\frac{y^\ell}{\ell!},
 \qquad y\ge0.
\]
The elementary bound
\begin{equation}\label{eq:ub-laguerre}
 0\le L_m(-y)
 \le\sum_{\ell\ge0}\frac{(my)^\ell}{(\ell!)^2}
 \le\cosh(2\sqrt{my})\le e^{2\sqrt{my}}
\end{equation}
follows from $\binom{2\ell}{\ell}\le4^\ell$.  In particular, if all
$b_j=\beta/2$, expansion of \eqref{eq:ub-volterra} over subsets gives
\begin{equation}\label{eq:ub-basis-majorants}
 \begin{split}
 |\varphi_j(t)|&\le\sqrt\beta e^{-\beta t/2}
                         L_{j-1}(-\beta t),\quad t\ge0,\\
 |\varphi_j(-x)|&\le\sqrt\beta e^{\beta x/2}
                         L_{j-1}(-\beta x),\quad x\ge0.
 \end{split}
\end{equation}
For each subset with $\ell$ factors, the convolution simplex has volume
$t^\ell/\ell!$ or $x^\ell/\ell!$.  Equality of the real parts supplies
the common exponential factor; taking absolute values removes the signs
of the Volterra factors.  There are $\binom{j-1}{\ell}$ such subsets.

\subsection{Bilateral near-edge extrapolation}

\begin{theorem}[Near-edge estimate]\label{thm:ub-edge}
Let $N\ge1$ and
$g(x)=\sum_{j=1}^N c_je^{i\mu_jx}$, with $c_j\in\mathbb C$ and
$\mu_j\in\mathbb R$.  For every $d\ge0$,
\begin{equation}\label{eq:ub-edge}
 \max\{|g(1+d)|,|g(-1-d)|\}
 \le \sqrt{\frac{250}{19}}\,N
       e^{5\sqrt2N\sqrt d}\|g\|_{L^2(-1,1)}.
\end{equation}
The bound has no frequency-separation hypothesis.
\end{theorem}

\begin{proof}
Zero coefficients can be removed and repeated frequencies merged.  The
result with the resulting number of terms implies the stated bound, so
assume that all $N$ frequencies are distinct.  The zero function is
immediate.  For any $N$-term real-frequency exponential sum $F$ and
$\beta>0$, put
\[
 I_\beta(F)=\int_0^\infty|F(t)|^2e^{-\beta t}\,dt,
 \qquad H(t)=e^{-\beta t/2}F(t).
\]
Expand $H$ in \eqref{eq:ub-basis}, with
$\zeta_j=\beta/2-i\mu_j$.  Cauchy--Schwarz and
\eqref{eq:ub-basis-majorants} yield
\[
 |F(t)|^2e^{-\beta t}
 \le I_\beta(F)\,\beta e^{-\beta t}
              \sum_{j=1}^N L_{j-1}(-\beta t)^2.
\]
Consequently, by \eqref{eq:ub-laguerre},
\begin{align*}
 I_\beta(F)^{-1}\int_{25N/\beta}^\infty
             |F(t)|^2e^{-\beta t}\,dt
 &\le N\int_{25N}^\infty e^{-y+4\sqrt{Ny}}\,dy\\
 &\le N\int_{25N}^\infty e^{-y/5}\,dy
   =5Ne^{-5N}<\frac1{20}.
\end{align*}
The second inequality uses $4\sqrt{Ny}\le4y/5$ for $y\ge25N$.
It follows that
\begin{equation}\label{eq:ub-finite-range}
 I_\beta(F)\le\frac{20}{19}
       \int_0^{25N/\beta}|F(t)|^2e^{-\beta t}\,dt.
\end{equation}
At a negative time, the second bound in
\eqref{eq:ub-basis-majorants} instead gives
\[
 |F(-x)|^2\le\beta\sum_{j=1}^N L_{j-1}(-\beta x)^2 I_\beta(F).
\]
The factor $e^{\beta x}$ cancels because
$H(-x)=e^{\beta x/2}F(-x)$.

Take $F(t)=g(1-t)$, $x=d$, and $\beta=25N/2$.  The truncation point
in \eqref{eq:ub-finite-range} is $2$, and
\[
 |g(1+d)|^2
 \le\frac{20}{19}\frac{25N}{2}
       N e^{10\sqrt2N\sqrt d}\int_{-1}^1|g(x)|^2\,dx.
\]
Taking square roots proves the right endpoint bound.  Apply the same
argument to $g(-x)$ for the left endpoint.
\end{proof}

An affine consequence will be useful.  If $J=[-e,e]$, $e>0$, and
$I=[-e(1+d),e(1+d)]$, then
\begin{equation}\label{eq:ub-affine-edge}
 \|g\|_{L^\infty(I)}^2
 \le\frac{250}{19e}N^2e^{10\sqrt2N\sqrt d}
                  \int_J|g|^2.
\end{equation}
The theorem proves this outside $J$ and at its endpoints.  For an interior
point $z$, apply the $d=0$ assertion to $[-e,z]$ and $[z,e]$.  If the
two lengths are $\ell_-$ and $\ell_+$, the resulting inequalities are
$\ell_\pm|g(z)|^2\le(500/19)N^2\int_{J_\pm}|g|^2$.
Adding them and using $\ell_-+\ell_+=2e$ proves
\eqref{eq:ub-affine-edge} there as well.

\subsection{A time-weighted estimate}

\begin{lemma}[Time-weighted decay]\label{lem:ub-delay}
Let $(E,\mu)$ be a measure space, and let
$\zeta_1,\ldots,\zeta_N\in\mathbb C$ satisfy
$1\le\operatorname{Re}\zeta_j\le4$.  For measurable finite
coefficients $c_j(z)$, define
\[
 f(t,z)=\sum_{j=1}^N c_j(z)e^{-\zeta_jt},\qquad
 k(t)=\int_E|f(t,z)|^3\,d\mu(z),\qquad
 K_T=\int_0^T k(t)\,dt.
\]
For every $T>0$ with $K_T<\infty$,
\begin{equation}\label{eq:ub-delay}
 \int_0^T t\,k(t)^{2/3}\,dt\le C N^{4/3}K_T^{2/3},
\end{equation}
where $C$ is absolute.
\end{lemma}

\begin{proof}
Merge coincident exponents, retaining the notation $N$ for their number.
In \eqref{eq:ub-volterra}, $1\le b_j\le4$.  Every convolution over
$\ell$ factors is bounded by $e^{-t}t^\ell/\ell!$, while $2b_j\le8$.
Therefore
\[
 |\varphi_j(t)|\le\sqrt8 e^{-t}L_{j-1}(-8t)
              \le\sqrt8 e^{-t+2\sqrt{8Nt}}.
\]
For almost every $z$, set $I_b(z)=\int_0^b|f(t,z)|^2\,dt$.
The coefficients are finite, so $I_\infty(z)<\infty$.
Expansion in the orthonormal basis gives
\[
 |f(t,z)|^2\le8N e^{-2t+4\sqrt{8Nt}}I_\infty(z).
\]
Take $S=128N$.  For $t\ge S$ the exponent is at most $-t$, so
\[
 \int_S^\infty|f(t,z)|^2\,dt
 \le8Ne^{-S}I_\infty(z)<\tfrac12 I_\infty(z).
\]
Thus $I_\infty\le2I_S$ and
$|f(t,z)|^2\le16Ne^{-t}I_S(z)$ for $t\ge S$.

Suppose $T\ge S$.  Minkowski's integral inequality in $L^{3/2}(E)$
and H\"older's inequality in time imply
\[
 \|I_S\|_{L^{3/2}(E)}
 \le\int_0^S k(u)^{2/3}\,du\le S^{1/3}K_T^{2/3}.
\]
Hence $k(t)^{2/3}\le16NS^{1/3}e^{-t}K_T^{2/3}$ for $t\ge S$,
and the tail contributes at most
$16NS^{1/3}(S+1)e^{-S}K_T^{2/3}\le C N^{4/3}K_T^{2/3}$.
On the initial interval, H\"older gives
\[
 \int_0^{\min(T,S)}t\,k(t)^{2/3}\,dt
 \le\left(\int_0^{\min(T,S)}t^3\,dt\right)^{1/3}K_T^{2/3}
 \le4^{-1/3}S^{4/3}K_T^{2/3}.
\]
If $T<S$ this is already the whole integral.  These two cases prove the
lemma.  Reinstating the original number of exponents only increases its
right-hand side.
\end{proof}

The use of $I_S$ with $S=128N$ is essential to this argument.  Replacing
it by $I_T$ would introduce $T^{1/3}$ through H\"older's inequality and
would not establish the stated uniform bound.

For later use, suppose the temporal decay rates of $G(s,z)$ have real
parts in $[\lambda,4\lambda]$, and let
$h(s)=\int_E|G(s,z)|^3\,d\mu$ and $H=\int_0^4h(s)\,ds$.
If $0\le\eta(s)\le C_\eta s$, rescale $t=\lambda s$ in
\eqref{eq:ub-delay}.  Since $K_{4\lambda}=\lambda H$, this gives
\begin{equation}\label{eq:ub-cutoff-time}
 \lambda\int_0^4\eta(s)h(s)^{2/3}\,ds
 \le C C_\eta N^{4/3}\lambda^{-1/3}H^{2/3}.
\end{equation}

\subsection{The physical model and the uniform theorem}

Write $\mathbb T=\mathbb R/(2\pi\mathbb Z)$.  Fix
$0<\delta_0<\delta$ and a function
$\vartheta\in C_c^\infty((-\delta,\delta))$ such that
$0\le\vartheta\le1$ and $\vartheta=1$ on $[-\delta_0,\delta_0]$.
Fix also $\eta\in C^1([0,4])$ with
$0\le\eta\le1$ and $\eta(0)=0$.
For $R>0$, $a\ge\max(1,\delta+2\delta_0)$ and $(a+\delta)R<\pi/2$, set
\begin{align}
 \xi_{a,R}(x)&=\vartheta(|x|/R-a),&
 \mathcal S^{\rm plat}_{a,R}
 &=\{x\in\mathbb T^3:\,||x|/R-a|\le\delta_0\},
 \label{eq:ub-geometry}\\
 F(t,x_2)&=\sum_{j=1}^q A_je^{-n_j^2t}
                   \cos(n_jx_2-\phi_j-n_jBt),
 \label{eq:ub-F}
\end{align}
where $q\ge1$, $1\le n_1<\cdots<n_q\le2n_1$ are integers,
$A_j\ge0$, and $\phi_j,B\in\mathbb R$.
The support of the cutoff lies in a Euclidean chart of the torus, and it
is extended by zero outside that chart.  The field $u=(0,B,F)$ is an exact
smooth Navier--Stokes solution with viscosity one and pressure zero:
its only nonzero evolution equation is
$F_t+BF_{x_2}-F_{x_2x_2}=0$.

Define the complete flux and the full plateau observation by
\begin{align}
 \mathcal T_{\boldsymbol n,R}
 &=\frac12\int_0^{4R^2}\!\int_{\mathbb T^3}
       \eta(t/R^2)B\,\partial_2\xi_{a,R}(x)F(t,x_2)^2\,dx\,dt,
 \label{eq:ub-T}\\
 M^{\rm plat}_{\boldsymbol n,R}
 &=\int_0^{4R^2}\!\int_{\mathcal S^{\rm plat}_{a,R}}
                         |F(t,x_2)|^3\,dx\,dt.
 \label{eq:ub-M}
\end{align}
The time cutoff occurs in the flux; the plateau observation uses the
entire unweighted clock.

\begin{theorem}[Complete-clock upper bound]\label{thm:ub-flux}
For the data above, put
\[
 \Delta_a=\frac{\delta+\delta_0}{a-\delta_0},\qquad
 \Phi_a=20\sqrt2q\sqrt{\Delta_a}.
\]
There is a constant $C$, depending only on
$\vartheta,\delta_0,\delta$ and $\|\eta'\|_\infty$, such that
\begin{equation}\label{eq:ub-main-refined}
 |\mathcal T_{\boldsymbol n,R}|
 \le C a^{2/3}R^{-1/3}
       \left(q^4+\frac{q^7}{a^2}+q^{13/3}\right)e^{\Phi_a}
       \bigl(M^{\rm plat}_{\boldsymbol n,R}\bigr)^{2/3}.
\end{equation}
In particular,
\begin{equation}\label{eq:ub-main}
 |\mathcal T_{\boldsymbol n,R}|
 \le C a^{2/3}R^{-1/3}q^7e^{\Phi_a}
       \bigl(M^{\rm plat}_{\boldsymbol n,R}\bigr)^{2/3}.
\end{equation}
The constant is independent of $q,R$, all frequencies, amplitudes,
phases, and $B$.
\end{theorem}

\begin{proof}
Use the scaled variables
\[
 \begin{gathered}
 s=t/R^2,\quad z=x_2/(aR),\quad G(s,z)=F(R^2s,aRz),\\
 \kappa_j=n_jaR,\quad \alpha=\kappa_1,\quad
 v=BR/a,\quad\lambda=(\alpha/a)^2.
 \end{gathered}
\]
Then $G_s+vG_z-a^{-2}G_{zz}=0$, and its $2q$ complex spatial
frequencies are $\pm\kappa_j$, with $\alpha\le\kappa_j\le2\alpha$.
Let
\[
 \begin{gathered}
 e_a=1-\delta_0/a,\quad E_a=[-e_a,e_a],\quad
 I_a=[-1-\delta/a,1+\delta/a],\\
 h(s)=\int_{E_a}|G(s,z)|^3\,dz,\quad H=\int_0^4h(s)\,ds.
 \end{gathered}
\]
The assumptions imply $e_a\ge1/2$, $\Delta_a\le1$, and uniformly
bounded upper and lower lengths for these intervals.
Equation \eqref{eq:ub-affine-edge} with $N=2q$, followed by H\"older,
gives
\begin{equation}\label{eq:ub-observation}
 \|G(s)\|_{L^\infty(I_a)}^2
 \le\frac{1000}{19e_a}q^2e^{\Phi_a}\int_{E_a}|G|^2
 \le Cq^2e^{\Phi_a}h(s)^{2/3}.
\end{equation}

The following two estimates from the literature are used in the remaining
argument.  Erd\'elyi's Markov inequality~\cite{Erdelyi2017} states, for
$f(x)=\sum_{r=1}^N c_re^{i\mu_rx}$ on $[0,1]$, that
\begin{equation}\label{eq:ub-markov-input}
 \|f'\|_\infty\le C
       \left(N^5+\sum_{r=1}^N\mu_r^2\right)^{1/2}\|f\|_\infty.
\end{equation}
The more precise expression in that paper is
$(1+\varepsilon_N)(108N^5+\sum_r\mu_r^2)^{1/2}$, with a uniformly
bounded factor $1+\varepsilon_N$.
K\'os's endpoint inequality~\cite{Kos2008}, recorded in
Erd\'elyi's paper as (1.2), is
\begin{equation}\label{eq:ub-kos-input}
 |f(0)|\le2N\|f\|_{L^2(0,1)},\qquad
 f(r)=\sum_{j=1}^N c_je^{\gamma_jr},\quad
 \operatorname{Re}\gamma_j\ge0.
\end{equation}

Map $I_a$ affinely onto $[0,1]$ in
\eqref{eq:ub-markov-input}.  Its length is bounded above and below,
and $\sum_j(\kappa_j^2+(-\kappa_j)^2)\le8q\alpha^2$.
Together with \eqref{eq:ub-observation}, this proves
\begin{equation}\label{eq:ub-derivative}
 \|G_z(s)\|_{L^\infty(I_a)}^2
 \le C(q^7+q^3\alpha^2)e^{\Phi_a}h(s)^{2/3}.
\end{equation}
For fixed $z$, the function $r\mapsto G(4-r,z)$, $0\le r\le1$,
has at most $2q$ exponential terms with exponents
$(\kappa_j/a)^2\pm iv\kappa_j$.  Applying
\eqref{eq:ub-kos-input}, cubing, and using H\"older on the unit
interval yields
\begin{equation}\label{eq:ub-endpoint}
 h(4)\le64q^3\int_3^4h(s)\,ds\le64q^3H,
 \qquad h(4)^{2/3}\le16q^2H^{2/3}.
\end{equation}

The radial geometry can be computed exactly.  At $y=x_2$, the
cross-sectional integral of the cutoff is
$2\pi\int_{|y|}^\infty r\vartheta(r/R-a)\,dr$.
Differentiating with respect to $y$ gives
\[
 \int_{\mathbb T^2}\partial_2\xi_{a,R}(x_1,y,x_3)\,dx_1\,dx_3
 =-2\pi y\vartheta(|y|/R-a).
\]
Thus, with
\[
 W_a(z)=-2\pi az\vartheta(a(|z|-1)),\qquad
 \Xi_a(z)=\int_{-\infty}^zW_a(r)\,dr,
\]
the exact flux conversion is
\begin{equation}\label{eq:ub-flux-conversion}
 \mathcal T_{\boldsymbol n,R}
 =\frac{a^2R^3}{2}\,v\int_0^4\eta(s)
                         \int_{\mathbb R}W_a(z)G(s,z)^2\,dz\,ds.
\end{equation}
The function $W_a$ is odd and has zero integral.  Its two support
intervals have length $O(a^{-1})$, and $|W_a|\le Ca$ there.  Consequently
\begin{equation}\label{eq:ub-kernel-bounds}
 \operatorname{supp}\Xi_a\subset I_a,\qquad
 \|\Xi_a\|_1+\|\Xi_a\|_\infty\le C,\qquad
 \|\Xi_a''\|_1=\|W_a'\|_1\le Ca.
\end{equation}
These bounds follow directly by differentiating $W_a$ on its support;
that support stays away from $z=0$.

Put $\mathcal E(s)=\int_{\mathbb R}\Xi_a(z)G(s,z)^2\,dz$.
Multiplying
$(G^2)_s+v(G^2)_z-a^{-2}(G^2)_{zz}=-2a^{-2}|G_z|^2$
by $\Xi_a$, and then integrating in space and time, gives
\begin{align}
 v\int_0^4\eta\int W_aG^2
 &=\eta(4)\mathcal E(4)-\int_0^4\eta'\mathcal E\,ds
 \nonumber\\
 &\quad-a^{-2}\int_0^4\eta\int\Xi_a''G^2
       +2a^{-2}\int_0^4\eta\int\Xi_a|G_z|^2.
 \label{eq:ub-energy}
\end{align}
There is no initial-face term because $\eta(0)=0$.
Using \eqref{eq:ub-observation}, \eqref{eq:ub-endpoint}, and
\eqref{eq:ub-kernel-bounds}, the first three terms on the right have
absolute sum at most
\begin{equation}\label{eq:ub-first-rows}
 C(q^4+q^2+a^{-1}q^2)e^{\Phi_a}H^{2/3}.
\end{equation}
Here $\int_0^4h^{2/3}\le4^{1/3}H^{2/3}$.
The last term, by \eqref{eq:ub-derivative}, is bounded by
\[
 Ce^{\Phi_a}\left[
       \frac{q^7}{a^2}\int_0^4\eta h^{2/3}\,ds
       +q^3\lambda\int_0^4\eta h^{2/3}\,ds\right].
\]
The first summand is at most $Cq^7a^{-2}H^{2/3}$.
If $\lambda\le1$, H\"older bounds the second by $Cq^3H^{2/3}$.
If $\lambda>1$, the temporal exponents of $G$ have real parts in
$[\lambda,4\lambda]$ and $\eta(s)\le\|\eta'\|_\infty s$.
Equation \eqref{eq:ub-cutoff-time}, with $N\le2q$, bounds that summand
by $Cq^{13/3}\lambda^{-1/3}H^{2/3}\le Cq^{13/3}H^{2/3}$.
Combining these estimates with \eqref{eq:ub-energy} gives
\begin{equation}\label{eq:ub-scaled-bound}
 \left|v\int_0^4\eta\int W_aG^2\right|
 \le C\left(q^4+\frac{q^7}{a^2}+q^{13/3}\right)
                  e^{\Phi_a}H^{2/3}.
\end{equation}

Finally, the plateau cross-sectional area at $x_2=aRz$, divided by
$R^2$, equals
\[
 \pi\left(
 [(a+\delta_0)^2-a^2z^2]_+
 -[(a-\delta_0)^2-a^2z^2]_+\right).
\]
For $z\in E_a$ it is exactly $4\pi a\delta_0$.  As
$dt\,dx_2=aR^3\,ds\,dz$, this proves
\begin{equation}\label{eq:ub-mass-conversion}
 M^{\rm plat}_{\boldsymbol n,R}\ge4\pi\delta_0a^2R^5H.
\end{equation}
Equations \eqref{eq:ub-flux-conversion}, \eqref{eq:ub-scaled-bound},
and \eqref{eq:ub-mass-conversion} prove
\eqref{eq:ub-main-refined}.  Since $a,q\ge1$,
$q^4+q^7/a^2+q^{13/3}\le3q^7$, proving \eqref{eq:ub-main}.
\end{proof}

\begin{remark}[A shrinking-gap regime]\label{rem:ub-window}
Let $a=pL$ with fixed $p>0$, and let
$\omega=\exp(-\gamma L^2)$ with fixed $\gamma>0$.  Set
$P=R^{-2}\omega M^{\rm plat}_{\boldsymbol n,R}$ and
$X=(\omega/R)[\mathcal T_{\boldsymbol n,R}]_+$.
Equation \eqref{eq:ub-main} becomes
\[
 X\le C a^{2/3}q^7\omega^{1/3}e^{20\sqrt2q\sqrt{\Delta_a}}P^{2/3}.
\]
If $q=q(L)=o(L^{5/2})$, then $\Delta_a=O(L^{-1})$ and
\[
 \lim_{L\to\infty}\frac1{L^2}
   \log\left(Ca^{2/3}q(L)^7\omega^{1/3}
                   e^{20\sqrt2q(L)\sqrt{\Delta_a}}\right)
 =-\frac\gamma3.
\]
This is the rate of the proved upper-bound coefficient.  It does not
assert that a particular flux attains that rate.  For the normalization
$\gamma=2/3969$ used in the research notes, the rate is $-2/11907$.
\end{remark}
